% (User input window)
% 2026-02-15 17:50:14（Asia/Singapore）

\paragraph{Geometric interpretation of \texorpdfstring{$\delta_q$}{deltaq}: spectral dimensional scales, quadratic channel maximal exponent, and sampling threshold}
Following the notation in Remark~\ref{rem:pom-collision-kernel-rh-phase}, \(\delta_q=\log(\Lambda_q^2/r_q)\) describes the competition intensity between the quadratic exponential scale \(\Lambda_q^{2m}\) of the non-Perron channel and the bulk scale \(r_q^{m}\).
This paragraph geometrizes the crossover exponent under \(\varphi\)-scaling and provides its auditable formulation as the ``strongest relative exponent from quadratic observables'' and the ``minimal sampling exponent threshold.''

\begin{definition}[bulk/bdry spectral dimensions and correlation codimension]\label{def:pom-deltaq-spectral-dimensions}
For each integer \(q\ge 2\), let \(r_q\) and \(\Lambda_q\) denote the Perron root and the maximal modulus of non-Perron roots of \(A_q\), respectively (Remark~\ref{rem:pom-collision-kernel-rh-phase}).
Under the geometric scaling \(\varepsilon_m=\varphi^{-m}\) (Proposition~\ref{prop:pom-renyi-dimension-spectrum}), define
\[
d_{\mathrm{bulk}}(q):=\frac{\log r_q}{\log\varphi},\qquad
d_{\mathrm{bdry}}(q):=\frac{\log \Lambda_q}{\log\varphi}.
\]
Define also the ``correlation boundary codimension''
\[
\mathrm{codim}_{\mathrm{corr}}(q):=\frac{\log r_q-2\log\Lambda_q}{\log\varphi}.
\]
\end{definition}

\begin{proposition}[Geometric identity for \texorpdfstring{$\delta_q$}{deltaq}]\label{prop:pom-deltaq-spectral-dimension-identity}
Under the notation of Definition~\ref{def:pom-deltaq-spectral-dimensions}, the following strict identities hold:
\[
\boxed{\;
\frac{\delta_q}{\log\varphi}
=2d_{\mathrm{bdry}}(q)-d_{\mathrm{bulk}}(q)
\;}
\qquad\text{and}\qquad
\boxed{\;
\mathrm{codim}_{\mathrm{corr}}(q)=-\frac{\delta_q}{\log\varphi}\; }.
\]
\end{proposition}
\begin{proof}
From \(\delta_q=2\log\Lambda_q-\log r_q\), dividing both sides by \(\log\varphi\) yields the first identity.
The second identity follows directly from the definition of \(\mathrm{codim}_{\mathrm{corr}}\) and the first identity.
\end{proof}

\begin{corollary}[Half-threshold criterion and crossover phase]\label{cor:pom-deltaq-half-threshold-criterion}
Under the same notation, the following conditions are strictly equivalent:
\begin{enumerate}
  \item \(\delta_q=0\);
  \item \(d_{\mathrm{bdry}}(q)=\tfrac12 d_{\mathrm{bulk}}(q)\);
  \item \(\beta^{\mathrm{RH}}_q=\tfrac12\) (Remark~\ref{rem:pom-collision-kernel-rh-phase}).
\end{enumerate}
Moreover, \(\delta_q>0\) (respectively \(\delta_q<0\)) holds if and only if
\(
d_{\mathrm{bdry}}(q)>\tfrac12 d_{\mathrm{bulk}}(q)
\)
(respectively \(<\)).
\end{corollary}
\begin{proof}
Direct substitution using Proposition~\ref{prop:pom-deltaq-spectral-dimension-identity} and \(\beta^{\mathrm{RH}}_q=\log\Lambda_q/\log r_q\).
\end{proof}

% (User input window)
% 2026-02-15 19:56:58（Asia/Singapore）
% According to an internal note dated 2026-02-14.

\begin{definition}[Spectral Condensation Index]\label{def:pom-spectral-condensation-index-gLambda}
For each integer \(q\ge 2\), following the notation \(r_q=\rho(A_q)\) and
\(
\Lambda_q=\max_{\mu\in\sigma(A_q)\setminus\{r_q\}}|\mu|
\)
(Remark \ref{rem:pom-collision-kernel-rh-phase}).
Define the spectral condensation index as
\[
\boxed{\;
g_q^{(\Lambda)}:=\frac{1}{q}\log\frac{r_q}{\Lambda_q}\;}
\qquad(\text{thus }\ \Lambda_q/r_q=\exp(-qg_q^{(\Lambda)})).
\]
\end{definition}

\begin{corollary}[Strict Identity of \texorpdfstring{$\delta_q$}{deltaq} and \texorpdfstring{$g_q^{(\Lambda)}$}{gLambda}]\label{cor:pom-deltaq-gLambda-identity}
Under the notation of Definition \ref{def:pom-spectral-condensation-index-gLambda}, the crossing exponent satisfies the strict identity
\[
\boxed{\;
\delta_q=\log\!\left(\frac{\Lambda_q^2}{r_q}\right)=\log r_q-2q\,g_q^{(\Lambda)}\;}
\]
thus \(\delta_q=0\) is equivalent to \(g_q^{(\Lambda)}=\frac{1}{2q}\log r_q\).
\end{corollary}
\begin{proof}
From \(g_q^{(\Lambda)}=\frac{1}{q}(\log r_q-\log\Lambda_q)\) we have \(\log\Lambda_q=\log r_q-qg_q^{(\Lambda)}\).
Substituting into \(\delta_q=2\log\Lambda_q-\log r_q\) yields the shown identity. QED.
\end{proof}



\begin{proposition}[Maximal relative exponent of centered quadratic observables]\label{prop:pom-deltaq-quadratic-observable-max-exponent}
Let \(A\) be a real square matrix with spectral radius \(r>0\), and let
\[
\Lambda:=\max_{\mu\in\sigma(A)\setminus\{r\}}|\mu|.
\]
Assume a Perron--Frobenius type decomposition
\[
A^m=r^m\Pi+R_m,\qquad \|R_m\|=O(\Lambda^m),
\]
where \(\Pi\) is the Perron projection.
For any linear functional \(\mathcal L(\cdot)\), define the linear observable \(Z_m:=\mathcal L(A^m)\), and let \(c:=\mathcal L(\Pi)\) and the centered observable \(\widetilde Z_m:=Z_m-c\,r^m\).
Then:
\begin{enumerate}
  \item For any \(\mathcal L\), \(|\widetilde Z_m|=O(\Lambda^m)\), hence
  \[
  \limsup_{m\to\infty}\frac{1}{m}\log|\widetilde Z_m|\le \log\Lambda.
  \]
  \item If \(\mathcal L\) does not vanish on the spectral projection reaching modulus \(\Lambda\), then there exists \(c_\ast>0\) such that along infinitely many \(m\), \(|\widetilde Z_m|\ge c_\ast \Lambda^m\), thus
  \[
  \limsup_{m\to\infty}\frac{1}{m}\log|\widetilde Z_m|=\log\Lambda.
  \]
  \item For any two pairs of linear functionals \(\mathcal L_1,\mathcal L_2\) and corresponding centered observables \(\widetilde Z^{(i)}_m\), their quadratic/bilinear observables
  \[
  Q_m:=(\widetilde Z_m)^2\quad\text{or}\quad Q_m:=\widetilde Z^{(1)}_m\,\widetilde Z^{(2)}_m
  \]
  satisfy \(|Q_m|=O(\Lambda^{2m})\), and when both functionals are non-degenerate,
  \[
  \limsup_{m\to\infty}\frac{1}{m}\log|Q_m|=2\log\Lambda.
  \]
  Thus the maximal exponent relative to the dominant scale \(r^m\) is
  \[
  \boxed{\;
  \limsup_{m\to\infty}\frac{1}{m}\log\frac{|Q_m|}{r^m}
  =\log\!\left(\frac{\Lambda^2}{r}\right)\; }.
  \]
\end{enumerate}
\end{proposition}
\begin{proof}
Substituting the decomposition \(A^m=r^m\Pi+R_m\) into \(Z_m=\mathcal L(A^m)\) yields
\(
\widetilde Z_m=\mathcal L(R_m)
\), hence (1) follows immediately from \(\|R_m\|=O(\Lambda^m)\).
(2) follows from the spectral projection expansion (or Jordan decomposition): if \(\mathcal L\) is nonzero on the spectral component at modulus \(\Lambda\), that component cannot cancel along an infinite subsequence, so \(\limsup\) achieves \(\log\Lambda\).
(3) follows immediately from (1) for the upper bound \(Q_m=O(\Lambda^{2m})\); in the non-degenerate case the lower bound is attained along the same subsequence, hence \(\limsup=2\log\Lambda\); the final expression is a rewrite of \(2\log\Lambda-\log r\) as \(\log(\Lambda^2/r)\).
\end{proof}

\begin{corollary}[Minimal sampling exponent threshold for crossover phase]\label{cor:pom-deltaq-sample-complexity-threshold}
Following the notation of Proposition~\ref{prop:pom-deltaq-quadratic-observable-max-exponent}, suppose there exists a sampleable random observable \(X_m\) satisfying
\[
\mathbb E[X_m]=c\,r^{m/2},\qquad \mathrm{Var}(X_m)=\Theta(\Lambda^{2m})
\]
(where \(\Lambda^{2m}\) represents the exponential scale of the centered quadratic channel).
Let \(\widehat X_{m,N}\) be the sample mean from \(N\) independent draws; then the relative mean squared error satisfies
\[
\frac{\mathbb E[(\widehat X_{m,N}-\mathbb E[X_m])^2]}{\mathbb E[X_m]^2}
=\frac{\mathrm{Var}(X_m)}{N\,\mathbb E[X_m]^2}
\asymp
\frac{1}{N}\exp\!\left(m\log\!\left(\frac{\Lambda^2}{r}\right)\right).
\]
Therefore, to keep the relative error \(o(1)\) as \(m\to\infty\), the necessary and (in constant-factor sense) sufficient sample size scaling is
\[
N\gtrsim \exp\!\left(m\log\!\left(\frac{\Lambda^2}{r}\right)\right).
\]
Applying this to \(A_q\), we obtain: \(\delta_q=\log(\Lambda_q^2/r_q)\) gives the sample complexity phase transition exponent from ``polynomial manageable'' to ``exponential threshold.''
\end{corollary}
\begin{proof}
Direct substitution using the variance formula for independent sample means
\(
\mathrm{Var}(\widehat X_{m,N})=\mathrm{Var}(X_m)/N
\).
\end{proof}

\begin{remark}[Minimal boundary-visible order and discrete auditability]\label{rem:pom-deltaq-minimal-visible-order-auditability}
From the sign changes
\(
\delta_2<0,\ \delta_3<0,\ \delta_4>0
\)
in Table~\ref{tab:fold_collision_kernel_rh_scan_q2_17}, one identifies \(q=4\) as the minimal crossover point in the replication-order direction (i.e., the minimal ``boundary-dominated order,'' Remark~\ref{rem:pom-collision-kernel-rh-phase}).
\par
Moreover, the triplet \((r_q,\Lambda_q,\delta_q)\) can in principle be audited completely from finite data without any external fitting: by the Hankel rank theorem (Theorem~\ref{thm:hankel-rank-minimal}) and the minimal realization functor \(\mathsf{MinH}\) (Remark~\ref{rem:pom-minh-functor}), the minimal recurrence polynomial \(P_q(\lambda)\) of the sequence \(S_q(m)\) is canonically determined up to isomorphism by finite Hankel data; \(r_q\) and \(\Lambda_q\) can then be extracted directly from the root moduli structure of \(P_q\), so \(\delta_q=\log(\Lambda_q^2/r_q)\) acquires a discrete reproducible spectral certificate (see Algorithm~\ref{alg:pom-recurrence-polynomial-spectral-extraction}).
\end{remark}


